%% LeafletFA -- Supplementary Material (Bioinformatics). Skeleton 2026-11-01; contents listed in BIOINFORMATICS_PLAN.md section 2.
%% Supplementary Note 1 (duplicate reads) is ported from GenomeBiology_Submission2026/manuscript/supplement/supplementary_notes.tex.
\documentclass[webpdf,traditional,large,namedate]{oup-authoring-template}
\usepackage[utf8]{inputenc}
\usepackage[T1]{fontenc}
\usepackage{graphicx}
\usepackage{float}
\usepackage{booktabs}
\usepackage{xcolor}
\usepackage{amsmath}
\newcommand{\societylogo}{}
\let\cite\citep
\InputIfFileExists{numbers}{}{}
\newcommand{\todo}[1]{\textcolor{red}{[TODO: #1]}}
\renewcommand{\thefigure}{S\arabic{figure}}
\renewcommand{\thetable}{S\arabic{table}}

\begin{document}
\journaltitle{Bioinformatics}
\title[LeafletFA]{Supplementary Material: LeafletFA}
\maketitle

\section*{Supplementary Methods}
\todo{Full model, inference settings, E-step, validation protocol (pre-registration and amendments 1--4).}

\section*{Supplementary Note 1: Duplicate reads and the effective number of molecules per event}
\label{suppnote:duplicates}
\todo{Port from supplementary\_notes.tex; numbers via the ledger.}

\section*{Supplementary Figures}
% One figure per item in BIOINFORMATICS_PLAN.md section 2 (S1--S8). Labels fig:S<n> are referenced from the main text as \ref{S-fig:S<n>}.
\begin{figure}[H]
\caption{\todo{Catalog coverage per dataset}}\label{fig:S1}
\end{figure}

\bibliographystyle{oup-abbrvnat}
\bibliography{references}
\end{document}
